\documentclass[11pt]{article}

% -------------------------------------------------
% Idioma y codificación
% -------------------------------------------------
\usepackage[utf8]{inputenc}
\usepackage[T1]{fontenc}
\usepackage[spanish,mexico]{babel}

% -------------------------------------------------
% Tipografía y composición
% -------------------------------------------------
\usepackage{lmodern}
\usepackage{microtype}
\usepackage{setspace}

% -------------------------------------------------
% Márgenes
% -------------------------------------------------
\usepackage[
    top=2.5cm,
    bottom=2.5cm,
    left=2.7cm,
    right=2.7cm
]{geometry}

% -------------------------------------------------
% Tablas, figuras y listas
% -------------------------------------------------
\usepackage{booktabs}
\usepackage{tabularx}
\usepackage{longtable}
\usepackage{array}
\usepackage{graphicx}
\usepackage{caption}
\usepackage{enumitem}

% -------------------------------------------------
% Matemáticas
% -------------------------------------------------
\usepackage{amsmath}
\usepackage{amssymb}

% -------------------------------------------------
% Referencias
% -------------------------------------------------
\usepackage{csquotes}
% Referencias integradas para compilar este archivo sin .bib ni Biber.
\usepackage[authoryear,round]{natbib}
\setlength{\bibhang}{1.5em}
\setlength{\bibsep}{0.6em}

% -------------------------------------------------
% Hipervínculos
% -------------------------------------------------
\usepackage[
    hidelinks
]{hyperref}

% -------------------------------------------------
% Formato general
% -------------------------------------------------
\setstretch{1.15}

\setlength{\parindent}{0pt}
\setlength{\parskip}{0.6em}

\setlist[itemize]{
    leftmargin=*,
    topsep=0.3em,
    itemsep=0.2em
}

\setlist[enumerate]{
    leftmargin=*,
    topsep=0.3em,
    itemsep=0.2em
}

\usepackage{placeins}
% -------------------------------------------------
% Documento
% -------------------------------------------------
\begin{document}

% -------------------------------------------------
% Portada
% -------------------------------------------------
\begin{titlepage}
    \centering

    \vspace*{2.0cm}

    {\large\bfseries
    Avance 2
    \par}

    \vspace{0.6cm}

    {\LARGE\bfseries
    \parbox{0.88\textwidth}{
        \centering
        Delimitación técnica, formulación y diseño experimental
        para la estimación de tiempos de viaje
    }
    \par}

    \vspace{1.2cm}

    {\large
    \parbox{0.88\textwidth}{
        \centering
        Planeación inteligente de rutas ante interrupciones mediante técnicas de optimización e inteligencia artificial
    }
    \par}

    \vspace{1.5cm}

    {\large\bfseries
    Equipo 42 -- Route Intelligence
    \par}

    \vspace{0.35cm}

    {\normalsize
    Maestría en Inteligencia Artificial Aplicada
    \par}

    \vspace{1.5cm}

    {\normalsize
    Luis Ángel Herrera Flores\\[0.25cm]
    Francisco Natanael Ortiz Martínez\\[0.25cm]
    Edgar Iván Rosas Gómez
    \par}

    \vspace{1.2cm}

    {\normalsize
    Asesor / Sponsor: Dr. Gerardo Jesús Camacho González
    \par}

    \vfill

    {\normalsize
    Octubre de 2026
    \par}

\end{titlepage}

% -------------------------------------------------
% Índice
% -------------------------------------------------
\tableofcontents
\newpage

\section{Introducción y delimitación del problema}

El Proyecto Integrador aborda la planeación y replaneación de rutas ante
interrupciones en problemas de ruteo de vehículos con ventanas de tiempo
(VRPTW, por \textit{Vehicle Routing Problem with Time Windows}). Como modelo
de referencia se considera el enfoque desarrollado por Berth y Trigos para el
problema de ruteo de vehículos con ventanas de tiempo ante interrupciones
(DVRPTW, por \textit{Disrupted Vehicle Routing Problem with Time Windows}), en
el que los tiempos asociados con los desplazamientos entre nodos constituyen
una entrada fundamental del modelo de optimización
\citep{berth_trigos_2025,berth_thesis_2025}. La calidad de esta información
resulta relevante tanto para construir la planeación inicial como para evaluar
y, en su caso, recuperar una operación cuando las condiciones de movilidad
cambian durante la ejecución.

Dentro de esta arquitectura se distinguen dos momentos de decisión. Antes de
iniciar la operación interesa anticipar las condiciones de movilidad esperadas
y construir una matriz de tiempos congruente con ellas; durante la ejecución,
una modificación no prevista puede volver subóptima o incluso inviable la
planeación original y exigir una actualización rápida de los tiempos relevantes
para los servicios pendientes. Esta distinción se desarrolla formalmente en la
Sección~\ref{sec:momentos}.

En el avance anterior se analizaron distintos puntos en los que técnicas de
inteligencia artificial o modelado probabilístico podrían complementar este
mecanismo sin sustituir necesariamente el procedimiento de optimización. La
estimación de los tiempos de recorrido se identificó como una alternativa
especialmente relevante, ya que permite concentrar el uso de estas técnicas en
mejorar uno de los parámetros que inciden directamente en las decisiones de
ruteo.

Esta línea de trabajo adquiere importancia porque los tiempos de viaje en una
red vial no son constantes. Presentan variaciones espaciales y temporales
asociadas con los patrones de movilidad, la congestión y cambios en las
condiciones de circulación. Una revisión sistemática reciente de la literatura
organiza los problemas de predicción de tráfico en tareas temporales,
espacio-temporales y origen--destino, y mantiene entre los retos centrales la
escasez de datos y la representación de dependencias espacio-temporales
\citep{du2025survey}. En esta línea, trabajos recientes han extendido la
estimación hacia ubicaciones no observadas a partir de información disponible
en otros puntos de la red \citep{lei2024conditional,zhou2025mogernn},
mientras que otros abordan de manera explícita la predicción de tiempos
origen--destino bajo información incompleta o dispersa
\citep{wang2025od}.

El presente avance mantiene los dos momentos como componentes centrales de la
arquitectura general: la planeación previa del siguiente turno y la
actualización durante la operación ante una disrupción. Para delimitar una
primera prueba de concepto, el diseño experimental se desarrolla con mayor
detalle para el segundo escenario, sin considerar la planeación previa como una
aplicación secundaria. Como hipótesis de solución propuesta por el equipo, se
estudiará si un modelo aprendido previamente a partir de información histórica
puede utilizar un número reducido de observaciones actuales de tiempo de
recorrido para estimar los tiempos requeridos por la replaneación. Estas
observaciones, denominadas consultas guía dentro del diseño experimental, no
tienen que coincidir con los clientes o depósitos de la instancia logística;
pueden seleccionarse de forma que aporten información útil sobre diferentes
zonas o corredores de la red.

En términos generales, el problema operacional de interés puede representarse
como

\[
K\text{ consultas actuales de tiempo de viaje}
\longrightarrow
\text{modelo previamente aprendido}
\longrightarrow
\widehat{T}^{\,\mathrm{rec}},
\]

donde \(\widehat{T}^{\,\mathrm{rec}}\) representa la matriz actualizada que
posteriormente puede utilizarse en el mecanismo de monitoreo y recuperación.
El propósito experimental será analizar cómo cambia la calidad de esta
estimación conforme varían tanto el número como la selección de las consultas
disponibles.

\section{Arquitectura metodológica y punto de integración con el DVRPTW}
\label{sec:arquitectura}

\subsection{Punto de integración con el modelo de referencia}

El enfoque planteado en este avance conserva como núcleo de decisión el
modelo de optimización utilizado para la planeación y recuperación de rutas.
La incorporación de técnicas de inteligencia artificial se sitúa
conceptualmente en una etapa previa, orientada a mejorar la información
temporal que posteriormente puede ser utilizada por el optimizador. De esta
manera, la formulación matemática del VRPTW y del DVRPTW se mantiene
inicialmente sin cambios, mientras que se define el punto de integración de
un componente destinado a estimar los tiempos asociados con los
desplazamientos entre los nodos de la instancia.

En el modelo de referencia, la matriz

\[
T=[t_{ij}]
\]

contiene los tiempos asociados con las transiciones entre los nodos
\(i,j\in N\). Estos valores intervienen tanto en la función objetivo como en
las restricciones que garantizan la consistencia temporal de las rutas
\citep{berth_thesis_2025,berth_trigos_2025}. Conforme a la convención
utilizada por Berth, \(t_{ij}\) incorpora además el tiempo de servicio
correspondiente al nodo de destino.

Para distinguir el componente que será objeto del diseño experimental, se
denotará mediante \(\tau_{ij}\) el tiempo requerido exclusivamente para el
desplazamiento entre los nodos \(i\) y \(j\). Por tanto,

\[
t_{ij}=\tau_{ij}+s_j,
\]

donde \(s_j\) representa el tiempo de servicio en el nodo \(j\). Bajo esta
notación, el componente asociado con las condiciones de movilidad es
\(\tau_{ij}\), mientras que \(s_j\) continúa siendo un parámetro propio de la
instancia logística.

En una implementación posterior, si un método de estimación produce un valor
\(\widehat{\tau}_{ij}\), el tiempo suministrado al modelo de optimización
podría construirse como

\[
\widehat{t}_{ij}=\widehat{\tau}_{ij}+s_j.
\]

El conjunto de estos valores permitiría formar una matriz estimada
\(\widehat{T}=[\widehat{t}_{ij}]\), compatible con la estructura de entrada
del modelo de referencia. La forma en que dicha matriz se obtiene depende del
momento de decisión considerado. En la planeación previa interesa anticipar
las condiciones esperadas para el periodo de operación, mientras que durante
la ejecución interesa actualizar con rapidez la información disponible ante
cambios en las condiciones de movilidad.

\subsection{Dos momentos para la estimación de tiempos de viaje}
\label{sec:momentos}

La necesidad de construir o actualizar la información temporal puede
presentarse en dos momentos distintos del proceso de ruteo. Aunque en ambos
casos el resultado requerido por el optimizador es una matriz compatible con
su formulación, los horizontes de decisión y la información disponible son
diferentes.

El primer momento corresponde a la \textit{planeación previa a la operación}.
En este caso se dispone de un horizonte de tiempo mayor antes del inicio del
servicio y el interés consiste en anticipar las condiciones de movilidad que
prevalecerán durante el periodo de ejecución. Para ello pueden aprovecharse
patrones históricos, variables temporales y, cuando se encuentren disponibles,
otras fuentes contextuales relevantes. Conceptualmente, este proceso puede
representarse como

\[
\begin{aligned}
\text{información histórica y contextual}
&\longrightarrow
\text{estimación de condiciones futuras} \\
&\longrightarrow
\widehat{T}^{\,\mathrm{plan}}
\longrightarrow
\text{VRPTW}.
\end{aligned}
\]

La matriz \(\widehat{T}^{\,\mathrm{plan}}\) representa una estimación de los
tiempos que podrían utilizarse para construir la planeación inicial. En este
escenario, el problema tiene un carácter principalmente predictivo y no se
identifica con el experimento de actualización rápida desarrollado más
adelante.

El segundo momento corresponde a la \textit{recuperación durante la
ejecución}. En este caso ya existe una planeación, una parte de los servicios
puede haber sido completada y una modificación inesperada de las condiciones
de movilidad puede requerir una nueva decisión. El tiempo disponible para
obtener información actualizada es entonces mucho menor.

Para este segundo escenario se propone separar explícitamente una fase de
aprendizaje previa y una fase de utilización durante la operación. A partir de
una base histórica \(\mathcal{D}_{\mathrm{hist}}\) se ajustará previamente un
modelo \(M_{\theta}\):

\[
\mathcal{D}_{\mathrm{hist}}
\longrightarrow
M_{\theta}.
\]

Cuando se presente una disrupción, el objetivo será utilizar únicamente un
conjunto reducido de consultas actuales de tiempo de recorrido como nueva
información sobre el estado de la red:

\[
\begin{aligned}
K\text{ consultas actuales}
&\longrightarrow
M_{\theta} \\
&\longrightarrow
\widehat{T}^{\,\mathrm{rec}}
\longrightarrow
\text{monitoreo y recuperación}.
\end{aligned}
\]

El histórico no se interpreta, por tanto, como una segunda fuente que deba
consultarse exhaustivamente al momento de la disrupción; su función principal
en este esquema es permitir que el modelo aprenda previamente relaciones
recurrentes entre las condiciones observables y los tiempos de viaje de la
red.

En la formulación de referencia, las modificaciones de los tiempos durante la
ejecución se representan mediante una matriz actualizada \(T'\), utilizada
para evaluar si las rutas continúan siendo factibles y, en caso necesario,
activar el proceso de recuperación
\citep{berth_thesis_2025,berth_trigos_2025}. Una matriz
\(\widehat{T}^{\,\mathrm{rec}}\) obtenida a partir de un número reducido de
consultas actuales podría proporcionar los valores requeridos por este
mecanismo sin necesidad de consultar individualmente todos los pares de la
instancia residual.

Los dos escenarios comparten la necesidad de transformar información sobre
movilidad en tiempos utilizables por el optimizador, pero no constituyen el
mismo problema. La planeación requiere anticipar condiciones futuras; la
recuperación requiere actualizar rápidamente la representación de las
condiciones presentes. Ambos forman parte del problema central del Proyecto.
Para mantener acotada la primera validación experimental, este avance desarrolla
con mayor detalle el segundo escenario, mientras que la estimación anticipada
para el siguiente turno se conserva como parte de la misma arquitectura y
deberá evaluarse en las etapas posteriores.

\section{Datos y entorno experimental}

\subsection{Origen de los datos y selección del conjunto experimental}

El estudio del problema planteado requiere disponer de una base histórica
suficientemente amplia para caracterizar la variación temporal y espacial de
las condiciones de movilidad. En una aplicación empresarial, esta información
podría provenir de registros acumulados durante la operación de una flota o de
servicios externos de información vial. Sin embargo, para construir un
experimento reproducible y evaluar inicialmente la viabilidad del enfoque se
consideraron conjuntos de datos públicos de tráfico con observaciones
repetidas sobre una misma red.

Una de las principales fuentes disponibles para este tipo de información es
el \textit{Performance Measurement System} (PeMS) del California Department of
Transportation (Caltrans). El sistema integra datos de detectores distribuidos
en la red de autopistas de California y conserva información histórica para su
análisis \citep{caltrans_pems}.

Para este proyecto se revisaron dos conjuntos utilizados como referencia en
investigaciones sobre predicción espacio-temporal de tráfico: PEMS-BAY y
METR-LA \citep{li2018dcrnn}. PEMS-BAY corresponde a una selección de
sensores ubicados en el área de la Bahía de San Francisco y se deriva de la
infraestructura PeMS. METR-LA constituye un conjunto de referencia diferente,
con mediciones de sensores del área de Los Ángeles. Ambos fueron utilizados
por Li et al. en el trabajo que introdujo la \textit{Diffusion Convolutional
Recurrent Neural Network} (DCRNN), un modelo de aprendizaje profundo para
predicción espacio-temporal de tráfico.

La variable principal disponible en estos conjuntos es la velocidad del
tráfico registrada en cada sensor a intervalos regulares de cinco minutos.
Por tanto, los datos no describen directamente viajes de vehículos de reparto
ni tiempos entre clientes. Su utilidad para este proyecto radica en que
proporcionan observaciones simultáneas y repetidas de diferentes puntos de
una red durante periodos prolongados, lo que permite caracterizar patrones
históricos y construir escenarios controlados de evaluación.

La versión de PEMS-BAY analizada contiene información de 325 sensores y
52\,116 instantes de observación correspondientes al primer semestre de 2017,
con una frecuencia nominal de cinco minutos. Cada registro puede representarse
como un vector de velocidades

\[
\mathbf{x}_t=
\begin{bmatrix}
v_1(t) & v_2(t) & \cdots & v_{325}(t)
\end{bmatrix}^{\mathsf T},
\]

de modo que el conjunto completo forma una serie temporal multivariada que
describe cómo evolucionan simultáneamente las velocidades observadas en
distintos puntos de la red.

Además de las series temporales de velocidad, la distribución utilizada
dispone de información complementaria para identificar los sensores y
representar relaciones espaciales entre ellos, incluyendo sus ubicaciones y
distancias. Esta información resulta relevante para el proyecto porque
permite vincular las observaciones temporales con una representación espacial
de la red y proporciona los elementos necesarios para construir el entorno
experimental descrito en la sección siguiente.

Como alternativa se examinó también METR-LA. Este conjunto contiene
información de 207 sensores, igualmente registrada con una frecuencia de
cinco minutos. La inspección preliminar de los archivos mostró una presencia
considerablemente mayor de valores iguales a cero y periodos prolongados de
indisponibilidad. En PEMS-BAY estos casos son sustancialmente menos frecuentes.
Por esta razón, PEMS-BAY se seleccionó como conjunto principal para la primera
etapa experimental, mientras que METR-LA se conserva como una posible fuente
para evaluar posteriormente la robustez de los métodos frente a información
incompleta.

\subsection{Función de PEMS-BAY dentro del experimento}

Aunque PEMS-BAY proporciona una descripción espacio-temporal rica de las
condiciones de tráfico, su estructura no coincide directamente con la
información requerida por una instancia de ruteo. Las observaciones
corresponden a velocidades registradas en sensores, mientras que el VRPTW y
el DVRPTW requieren tiempos de desplazamiento entre los nodos logísticos de
la instancia. Por ello, el conjunto de datos se utilizará como laboratorio
para construir y evaluar el mecanismo de estimación propuesto.

En un instante \(t\), el estado observado puede representarse mediante el
vector

\[
\mathbf{x}_t=
\begin{bmatrix}
v_1(t) & v_2(t) & \cdots & v_{n_s}(t)
\end{bmatrix}^{\mathsf T},
\]

donde \(v_r(t)\) representa la velocidad registrada por el sensor \(r\) y
\(n_s\) corresponde al número de sensores considerados. La secuencia histórica
de estos vectores constituye la base para construir, mediante el procedimiento
definido en la Sección~\ref{sec:tiempos}, series de tiempos de desplazamiento de referencia y
observaciones históricas derivadas. Estas relaciones podrán utilizarse
posteriormente para ajustar o calibrar el modelo \(M_{\theta}\), sin requerir
que su formulación opere directamente sobre las velocidades de los sensores.

PEMS-BAY cumple además una segunda función: para cada instante de evaluación,
la disponibilidad simultánea de un conjunto amplio de mediciones permite
construir una referencia experimental de las condiciones de movilidad y de
los tiempos de desplazamiento derivados de ellas. Esta referencia se
utilizará exclusivamente para evaluar las estimaciones producidas por el
modelo.

Finalmente, el conjunto de datos permitirá simular el escenario de acceso
limitado a información actual. Una vez construida la referencia completa para
un instante \(t\), se seleccionará un número reducido de pares de puntos y se
calcularán únicamente sus tiempos de recorrido como si correspondieran a las
pocas consultas que podrían realizarse durante una disrupción. El modelo
recibirá estas consultas como información actual, mientras que el resto de la
referencia permanecerá oculto para fines de estimación.

De esta manera, PEMS-BAY desempeña tres funciones diferenciadas dentro del
experimento: proporciona la base histórica a partir de la cual se construirán
las observaciones utilizadas para el aprendizaje previo, permite obtener una
referencia controlada para evaluar los resultados y hace posible simular un
número limitado de consultas actuales sin depender, en esta etapa, de un
servicio externo en tiempo real.

La referencia construida a partir de PEMS-BAY se utilizará como una
\textit{referencia experimental derivada} y no como una observación directa de
los tiempos origen--destino reales de vehículos en la red vial. Sus valores se
obtienen a partir de las mediciones disponibles y de las hipótesis utilizadas
para transformar velocidades y relaciones espaciales en tiempos de
desplazamiento. Por tanto, su función es proporcionar un entorno controlado y
reproducible para comparar estimaciones, no representar una ``verdad física''
exacta de todos los tiempos de recorrido.

Finalmente, en esta etapa se parte del supuesto de que existe una base
histórica suficientemente amplia para caracterizar los patrones de movilidad
de la red. La construcción progresiva de dicha base a partir de la operación
cotidiana de una flota, así como su validación con información proveniente de
una empresa real, se consideran extensiones posteriores del proyecto.


\subsection{Criterios de preparación e ingeniería de características}
\label{sec:preparacion}

En el contexto de CRISP-ML, la preparación de los datos debe distinguir las
características de entrada de los tiempos objetivo que se desea estimar. Las
consultas guía constituyen información observable durante la operación,
mientras que los tiempos de referencia de los pares logísticos se reservan
como objetivos de aprendizaje o evaluación. Estos últimos no deberán
incorporarse como entradas actuales del estimador.

Como complemento del diseño experimental, se proponen las siguientes
características candidatas. Su inclusión definitiva deberá justificarse con
pruebas sobre los periodos de entrenamiento y validación; esta propuesta no
constituye evidencia de transformaciones ya ejecutadas.

\begin{table}[ht]
\centering
\small
\caption{Características candidatas para el componente de estimación}
\label{tab:caracteristicas}
\renewcommand{\arraystretch}{1.15}
\begin{tabularx}{\textwidth}{@{}p{0.23\textwidth}X X@{}}
\toprule
\textbf{Característica} & \textbf{Construcción} & \textbf{Justificación} \\
\midrule
Tiempo consultado & \(q_k(t)\), expresado en una unidad temporal común. &
Representa la información actual disponible para cada par guía. \\
Desviación histórica & \(q_k(t)-b_k(h(t))\), donde \(b_k\) es la mediana
histórica de la consulta para una franja horaria comparable. &
Distingue el cambio actual del patrón recurrente del recorrido. \\
Razón respecto al histórico & \(q_k(t)/b_k(h(t))\), únicamente cuando la
referencia histórica sea positiva y válida. &
Permite comparar cambios relativos entre recorridos de distinta duración. \\
Hora cíclica & \(\sin(2\pi h(t)/24)\) y \(\cos(2\pi h(t)/24)\), con
\(h(t)\) expresada en horas. &
Representa la continuidad entre el final y el inicio del día sin imponer un
salto artificial entre las 23:59 y las 00:00 horas. \\
\bottomrule
\end{tabularx}
\end{table}
\FloatBarrier

Las referencias históricas \(b_k\) se calcularán exclusivamente con datos de
entrenamiento. Antes de construir los tiempos se documentarán los valores
faltantes, las velocidades cero y las observaciones no válidas; no se
interpretará automáticamente un cero como dato ausente ni se calculará
\(d_e/v_e(t)\) con un denominador nulo. La regla de tratamiento y su impacto
sobre los registros disponibles deberán quedar registrados.

Para métodos sensibles a la escala se evaluará la estandarización de las
características mediante \(z=(x-\mu_{\mathrm{train}})/\sigma_{\mathrm{train}}\).
Las variables con varianza nula se tratarán antes de esta operación. La
transformación logarítmica se considerará únicamente para variables positivas
cuando su distribución y las pruebas de validación la justifiquen; no se
presupondrá que todas las entradas deban seguir una distribución normal. Las
estadísticas y transformaciones ajustadas en entrenamiento se conservarán
para validación, prueba y utilización posterior.

La selección considerará filtros de varianza y correlación para detectar
características poco variables o redundantes. Sus umbrales deberán
justificarse con los resultados de validación. Como alternativa de extracción
se evaluará PCA sobre las entradas estandarizadas, ajustado únicamente con
entrenamiento; el número de componentes deberá fundamentarse en la varianza
explicada y en la calidad de estimación. Se registrarán las dimensiones antes
y después, así como el efecto sobre el tiempo de procesamiento. La selección
espacial de consultas complementa estas pruebas, pero no sustituye el análisis
de las características construidas.

\section{Formulación del entorno experimental y actualización mediante consultas guía}
\label{sec:entorno}

\subsection{Construcción de los tiempos de referencia}
\label{sec:tiempos}

Para relacionar las observaciones de PEMS-BAY con el problema de estimación
de tiempos de viaje, se propone representar el entorno experimental mediante
un grafo

\[
G=(V,E),
\]

donde \(V\) contiene los puntos utilizados para representar espacialmente la
red y \(E\) las relaciones consideradas entre ellos. Esta representación se
utiliza con fines experimentales y no implica que cada relación de los
archivos complementarios de PEMS-BAY deba interpretarse necesariamente como
un tramo vial físico directo.

Para cada conexión \(e\in E\) se asociará una distancia \(d_e\) y se definirá
una velocidad \(v_e(t)\) a partir de las mediciones disponibles en el instante
\(t\). A partir de estos valores se obtendrá un peso temporal para la
conexión,

\[
w_e(t)=\frac{d_e}{v_e(t)},
\]

una vez expresadas ambas variables en unidades compatibles. La regla
específica para asignar una velocidad a cada conexión deberá establecerse y
documentarse como una hipótesis del diseño experimental, ya que PEMS-BAY
proporciona mediciones en sensores y no tiempos de recorrido directamente
asociados con todas las conexiones de una red vial.

Sean \(a,b\in V\) dos puntos de la red experimental y
\(\mathcal{P}_{ab}\) el conjunto de caminos factibles que los conectan. El
tiempo de desplazamiento de referencia entre ambos se definirá mediante el
camino de menor tiempo,

\[
\tau_{ab}^{\mathrm{ref}}(t)
=
\min_{p\in\mathcal{P}_{ab}}
\sum_{e\in p} w_e(t).
\]

Esta formulación permite abstraer los nodos y conexiones intermedios utilizados
para construir cada tiempo origen--destino. El recorrido interno puede cambiar
entre distintos instantes conforme cambien los pesos temporales de las
conexiones, pero para las etapas posteriores el resultado relevante es el
tiempo agregado \(\tau_{ab}^{\mathrm{ref}}(t)\).

Dentro de \(V\) se definirá un subconjunto \(C\subseteq V\) correspondiente a
los nodos logísticos para los que el modelo de optimización requiere tiempos
de desplazamiento. Para estos pares, los valores

\[
\left\{
\tau_{ij}^{\mathrm{ref}}(t)
\right\}_{i,j\in C}
\]

constituyen la referencia contra la cual se evaluarán posteriormente las
estimaciones \(\widehat{\tau}_{ij}(t)\). La construcción de tiempos para pares
de puntos que no pertenecen a \(C\) se conserva porque dichos pares pueden
utilizarse como consultas guía para obtener información sobre el estado actual
de la red.

Para una primera implementación se considerará que los pesos temporales
\(w_e(t)\) permanecen aproximadamente constantes dentro de cada intervalo de
cinco minutos del conjunto de datos. Bajo este supuesto, todos los pesos
utilizados para calcular un camino corresponden al mismo estado temporal de
la red. Una extensión posterior podría considerar un problema de camino
mínimo dependiente del tiempo, en el que el costo de cada conexión se evalúe
de acuerdo con el instante en que el vehículo ingresa a ella.

\subsection{Consultas guía como observaciones del estado actual}

Como hipótesis de solución propuesta por el equipo, el mecanismo de
actualización investigará bajo qué condiciones un número reducido de consultas
actuales permite estimar los tiempos de viaje de los pares logísticos no
observados con precisión suficiente para la replaneación. Esta formulación no
se asume como requisito del modelo de referencia, sino como una alternativa
para reducir la cantidad de información actual que tendría que obtenerse antes
de reconstruir la matriz requerida por el optimizador. Se plantea seleccionar
previamente los pares consultados por su capacidad potencial para aportar
información sobre diferentes zonas o corredores de la red. Esta capacidad
deberá comprobarse experimentalmente, especialmente ante cambios localizados
que las consultas podrían no detectar.

Sea

\[
\mathcal{Q}_K=
\left\{
(a_k,b_k)
\right\}_{k=1}^{K}
\]

el conjunto de \(K\) consultas guía. Los extremos \(a_k,b_k\in V\) no tienen
que pertenecer al conjunto de nodos logísticos \(C\). Para cada consulta, la
información disponible durante la operación consiste únicamente en el tiempo
de recorrido observado o consultado entre ambos puntos,

\[
q_k(t)=\tau_{a_kb_k}(t).
\]

En el entorno experimental con PEMS-BAY, cada valor \(q_k(t)\) se simulará
mediante el tiempo de referencia correspondiente
\(\tau_{a_kb_k}^{\mathrm{ref}}(t)\), calculado a partir del estado completo del
instante de evaluación. Esta equivalencia se utiliza únicamente para construir
un escenario controlado en el que el modelo observa las consultas seleccionadas
mientras los tiempos objetivo permanecen ocultos.

En una aplicación real, estos valores podrían obtenerse mediante un servicio
externo de mapas e información de tráfico, como Google Maps, o a partir de
otras fuentes operativas. El método propuesto no requiere conocer
explícitamente la secuencia de calles o conexiones utilizada por el servicio
para producir cada tiempo: la observación se trata como un tiempo agregado
entre dos puntos conocidos.

Las consultas guía pueden interpretarse conceptualmente como \textit{sensores
virtuales agregados}: cada una resume en un único tiempo de viaje las
condiciones de movilidad que afectan un recorrido de la red. Esta expresión se
utiliza únicamente como una analogía conceptual; en el documento se mantiene
el término \textit{consulta guía} para evitar confundirlas con sensores físicos
o con una categoría establecida de la literatura.

El modelo \(M_{\theta}\), ajustado previamente con información histórica,
recibirá durante la operación el conjunto de observaciones actuales

\[
\mathbf{q}_t=
\begin{bmatrix}
q_1(t) & q_2(t) & \cdots & q_K(t)
\end{bmatrix}^{\mathsf T}
\]

y producirá estimaciones de los tiempos de desplazamiento requeridos entre los
nodos logísticos,

\[
\left\{
\widehat{\tau}_{ij}(t)
\right\}_{i,j\in C}.
\]

Por tanto, el problema de recuperación estudiado experimentalmente puede
representarse como

\[
\mathbf{q}_t
\xrightarrow{\;M_{\theta}\;}
\left\{
\widehat{\tau}_{ij}(t)
\right\}_{i,j\in C},
\]

donde el conocimiento histórico necesario para establecer las relaciones
entre distintas condiciones de movilidad se encuentra incorporado previamente
en los parámetros del modelo.

Esta formulación conserva la separación entre la capa de movilidad y el modelo
de optimización. Una vez obtenidos los valores
\(\widehat{\tau}_{ij}(t)\), la incorporación de los tiempos de servicio
\(s_j\) y la construcción de \(\widehat{T}^{\,\mathrm{rec}}\) se realizará de
acuerdo con la relación establecida en la Sección~\ref{sec:arquitectura}.

\subsection{Diseño de escenarios y selección de consultas}
\label{sec:escenarios}

El diseño experimental estudiará dos dimensiones relacionadas: cuántas
consultas actuales se requieren y qué pares conviene consultar. Para un mismo
instante de referencia se evaluarán distintos valores de \(K\), manteniendo
ocultos para el modelo los tiempos correspondientes a los pares logísticos
que se desea estimar.

En una primera etapa, la selección de consultas deberá permitir comparar el
valor de una cobertura deliberada de la red frente a una selección sin
criterio informativo. Por ello se considerará al menos una estrategia de
referencia basada en selección aleatoria y una estrategia estructurada que
busque distribuir las consultas entre distintas zonas o corredores y reducir
la redundancia espacial. El objetivo inicial no es resolver un problema
óptimo de localización de consultas, sino comprobar si la elección de los
pares observados modifica de manera material la calidad de la estimación.

De forma esquemática, la pregunta experimental puede expresarse como

\[
(K,\,\mathcal{Q}_K)
\longrightarrow
\text{calidad de}
\left\{
\widehat{\tau}_{ij}(t)
\right\}_{i,j\in C}.
\]

En una extensión posterior, la selección de consultas podría formalizarse
mediante criterios de información, cobertura u observabilidad, de modo que un
número reducido de observaciones proporcione la mayor cantidad posible de
información sobre el estado relevante de la red. Esta extensión no se
considera un requisito para la primera prueba de concepto.

Para cada instante de evaluación, PEMS-BAY permitirá construir tanto los
tiempos de referencia de los pares logísticos como los valores que simulan
las consultas guía. Al modelo se entregarán únicamente estas últimas, mientras
que los valores de referencia se conservarán para medir el error de las
estimaciones. La comparación se repetirá sobre múltiples instantes y
condiciones de circulación con el fin de evitar que los resultados dependan
de un único estado de la red.

\section{Compatibilidad metodológica y criterios de evaluación}

La estimación y predicción de información de tráfico comprende tareas con
funciones diferentes dentro de un sistema de movilidad. Una revisión
sistemática reciente distingue, entre otras, la predicción de series
temporales, la predicción espacio-temporal y los problemas origen--destino
\citep{du2025survey}. Esta clasificación resulta útil para el presente
proyecto porque algunos enfoques buscan anticipar condiciones futuras, otros
estiman condiciones actuales utilizando información parcial y otros producen
tiempos o variables asociadas con pares origen--destino. Estas funciones son
complementarias, pero requieren información distinta y no deben considerarse
intercambiables.

Conforme a la delimitación establecida en la Sección~\ref{sec:momentos}, el
Proyecto considera como centrales tanto la estimación anticipada para la
planeación previa como la actualización rápida durante la operación. En esta
primera validación, el análisis metodológico se desarrolla con mayor detalle
para el segundo escenario. En este contexto, la compatibilidad de una
alternativa depende de su capacidad para aprovechar información histórica,
incorporar observaciones actuales parciales y producir, directa o
indirectamente, estimaciones de los tiempos requeridos entre los nodos
logísticos. De acuerdo con la formulación
de la Sección~\ref{sec:entorno}, la transformación de interés conserva la forma

\[
\mathbf{q}_t
\xrightarrow{\;M_{\theta}\;}
\left\{
\widehat{\tau}_{ij}(t)
\right\}_{i,j\in C}.
\]


\subsection{Compatibilidad de enfoques con la actualización rápida}

Entre los antecedentes relacionados con esta formulación se encuentran los
modelos multivariados que explotan dependencias espacio-temporales aprendidas
a partir de información histórica. Jenelius y Koutsopoulos proponen un modelo
basado en \textit{Probabilistic Principal Component Analysis} (PPCA) que
representa correlaciones entre tiempos de viaje de una red y permite realizar
predicciones condicionadas a observaciones recientes aun cuando parte de la
información se encuentra ausente \citep{jenelius2018urban}. Cebecauer,
Jenelius y Burghout desarrollan posteriormente un marco de predicción de red
orientado al uso de datos de vehículos sonda escasos y valores faltantes,
combinando PPCA con información local
\citep{cebecauer2018integrated}.

La relevancia de estos trabajos para el problema planteado radica en que
muestran que las condiciones observadas en diferentes partes de una red pueden
presentar una estructura conjunta aprovechable para inferir información no
observada directamente. Sin embargo, sus formulaciones originales no utilizan
el mismo tipo de consultas origen--destino agregadas definidas en este
proyecto. Por tanto, su aplicación al esquema de consultas guía requeriría una
adaptación y una validación experimental; los resultados reportados en dichos
trabajos no permiten concluir por sí mismos cuántas consultas serían
suficientes ni cuáles deberían seleccionarse.

Investigaciones más recientes muestran que la inferencia a partir de cobertura
parcial de una red continúa siendo un problema activo. Lei, Xu y Wang proponen
un modelo probabilístico de difusión condicionada para estimar estados de
tráfico en ubicaciones sin sensores a partir de las dependencias
espacio-temporales con ubicaciones observadas \citep{lei2024conditional}.
Por su parte, Zhou et al. desarrollan un modelo inductivo basado en grafos para
predecir estados en ubicaciones no observadas y analizan su desempeño bajo
distintos niveles de disponibilidad de sensores \citep{zhou2025mogernn}.
Estos enfoques constituyen evidencia contemporánea de que observaciones
parciales pueden contener información útil sobre zonas no observadas, aunque
trabajan con estados asociados con nodos de una red explícita y no con el
esquema de consultas agregadas planteado aquí.

Otros enfoques probabilísticos proporcionan antecedentes complementarios.
Hofleitner et al. utilizan una red bayesiana dinámica para representar
dependencias espacio-temporales y actualizar la estimación de las condiciones
de tráfico a partir de observaciones dispersas de vehículos sonda
\citep{hofleitner2012learning}. Este esquema resulta conceptualmente
compatible con la combinación de conocimiento histórico y evidencia actual,
aunque su formulación mantiene estados asociados con los enlaces de la red y
requiere un procedimiento de inferencia más complejo que la representación
agregada utilizada en este avance.

Zhan et al. estudian un problema distinto pero especialmente relacionado con
la disponibilidad de información parcial. Su modelo estima tiempos medios de
enlaces utilizando datos de viajes que contienen origen, destino, duración y
distancia, sin observar directamente la trayectoria seguida; los caminos se
tratan como variables latentes dentro de un modelo de mezcla bayesiano
\citep{zhan2016bayesian}. Este antecedente muestra que observaciones
agregadas de viajes pueden aportar información sobre condiciones no observadas
directamente en la red, aunque su metodología requiere inferencia de caminos y
un volumen de observaciones distinto del escenario de pocas consultas
considerado aquí.

Un trabajo reciente se aproxima al problema desde la predicción directa del
tiempo origen--destino. Wang et al. combinan reconstrucción de datos faltantes,
representación adaptativa del grafo y modelado temporal para estimar tiempos OD
en presencia de información vial dispersa \citep{wang2025od}. Aunque su
objetivo de salida es cercano al de este proyecto, el método requiere una
representación más rica de la red y no aborda específicamente la reconstrucción
de una matriz a partir de un número reducido de consultas actuales.

También resulta relevante la cuestión del costo de adquirir tiempos OD desde
servicios externos. Shah y Piantanakulchai estudian un esquema de aprendizaje
activo que selecciona de manera iterativa qué tiempos consultar a una API para
construir un conjunto de entrenamiento con menor esfuerzo de etiquetado
\citep{shah2026active}. Esta propuesta respalda la utilidad de seleccionar
consultas cuando obtener una matriz completa es costoso, pero difiere del
problema operacional considerado aquí: en ese trabajo las consultas se utilizan
para incorporar nuevas etiquetas y reentrenar el modelo, mientras que en el
presente diseño las consultas se obtienen después del entrenamiento y se usan
como evidencia actual para producir \(
\widehat{T}^{\,\mathrm{rec}}
\) sin reentrenamiento en el momento de la disrupción.

También existen métodos cuya función se encuentra en otras etapas del proceso.
Modelos espacio-temporales sobre grafos, como DCRNN, están orientados
principalmente a predecir la evolución futura de variables de tráfico a partir
de series históricas distribuidas sobre una red
\citep{li2018dcrnn}. Por ello, resultan más directamente relacionados con
el problema de planeación previa que con la actualización mediante unas pocas
consultas actuales. En sentido diferente, Gentile et al. estudian la
generación eficiente de tiempos punto a punto cuando la información dinámica
de la red ya se encuentra disponible \citep{gentile2020fast}. Su propuesta
es, por tanto, complementaria a la estimación del estado actual, pero no
sustituye el mecanismo que debe inferir información a partir de observaciones
limitadas.

En conjunto, estos antecedentes respaldan la viabilidad de aprovechar
dependencias históricas y observaciones parciales para estimar información de
movilidad, pero también muestran que las alternativas disponibles resuelven
subproblemas diferentes. En esta etapa no se selecciona un método definitivo.
La comparación e implementación de modelos candidatos deberá realizarse
posteriormente sobre la formulación y los escenarios definidos en este
avance, manteniendo como requisito que la información nueva necesaria durante
la recuperación se limite al conjunto de consultas actuales disponible.


\subsection{Criterios de evaluación del componente de estimación}

La evaluación del mecanismo propuesto debe determinar si las estimaciones
obtenidas a partir de un conjunto reducido de consultas representan con
suficiente precisión los tiempos de desplazamiento que permanecerían ocultos
para el modelo. Para ello, los valores estimados se compararán con los tiempos
de referencia construidos mediante el procedimiento de la Sección~\ref{sec:entorno}.

Sea \(\mathcal{T}_{\mathrm{test}}\) el conjunto de instantes reservados para
evaluación. Para evitar que los resultados reflejen información utilizada
durante el ajuste del modelo, los periodos destinados a entrenamiento y prueba
deberán mantenerse separados respetando el orden temporal de los datos. En la
primera evaluación se considerará la matriz completa de pares ordenados de
nodos \(i,j\in C\), con \(i\neq j\). Bajo este criterio, el error absoluto
medio de los tiempos de desplazamiento podrá calcularse como

\[
\mathrm{MAE}_{\tau}
=
\frac{1}
{|\mathcal{T}_{\mathrm{test}}|\,|C|(|C|-1)}
\sum_{t\in\mathcal{T}_{\mathrm{test}}}
\sum_{\substack{i,j\in C\\ i\neq j}}
\left|
\widehat{\tau}_{ij}(t)
-
\tau_{ij}^{\mathrm{ref}}(t)
\right|.
\]

De manera complementaria se utilizará la raíz del error cuadrático medio,

\[
\mathrm{RMSE}_{\tau}
=
\sqrt{
\frac{1}
{|\mathcal{T}_{\mathrm{test}}|\,|C|(|C|-1)}
\sum_{t\in\mathcal{T}_{\mathrm{test}}}
\sum_{\substack{i,j\in C\\ i\neq j}}
\left(
\widehat{\tau}_{ij}(t)
-
\tau_{ij}^{\mathrm{ref}}(t)
\right)^2
}.
\]

El MAE permitirá interpretar directamente la magnitud promedio de los errores
en unidades de tiempo, mientras que el RMSE otorgará mayor peso a errores de
gran magnitud que podrían resultar especialmente relevantes para una decisión
de replaneación. Las métricas se aplicarán sobre los tiempos de desplazamiento
\(\tau_{ij}\), ya que los tiempos de servicio \(s_j\) permanecen como
parámetros de la instancia logística y se incorporan posteriormente al
construir \(\widehat{T}^{\,\mathrm{rec}}\).

El número de consultas \(K\) no se considerará una métrica de error, sino una
variable del diseño experimental. Para cada alternativa se analizará cómo
cambia la precisión al modificar \(K\) y la composición de
\(\mathcal{Q}_K\). Esto permitirá estudiar por separado el efecto de disponer
de más observaciones y el efecto de seleccionar pares potencialmente más
informativos. Conforme al diseño establecido en la Sección~\ref{sec:escenarios}, las
comparaciones deberán realizarse sobre los mismos instantes y tiempos objetivo
para evitar que diferencias en las condiciones de circulación se confundan
con diferencias atribuibles a las consultas seleccionadas.

Además de la precisión, se registrará el tiempo requerido por el modelo para
producir las estimaciones una vez disponibles las consultas guía. Este tiempo
corresponde a la etapa de inferencia y debe distinguirse del tiempo que, en una
aplicación real, pudiera requerir la obtención de información desde un
servicio externo. El entorno experimental con PEMS-BAY permite controlar el
primero y estudiar el efecto de \(K\), pero no pretende reproducir la latencia
de una plataforma comercial de información vial.

Para determinar si las observaciones actuales aportan información adicional
respecto de los patrones recurrentes de la red, los resultados se
compararán también con una referencia histórica sencilla construida sin
utilizar las consultas del instante de evaluación, por ejemplo mediante
valores característicos para periodos temporales comparables. Esta referencia
permitirá distinguir la mejora atribuible a la incorporación de información
actual de aquella que podría obtenerse únicamente a partir del comportamiento
histórico. Esta comparación incluirá el caso sin consultas actuales
(\(K=0\)). La elección de características, transformaciones y parámetros
utilizará un periodo de validación separado del conjunto final de prueba.
Para cada combinación de \(K\) y \(\mathcal{Q}_K\), se documentará qué
entradas recibe el modelo y si requiere un ajuste independiente; los pares
consultados deberán conservar una identidad y un orden explícitos.

La evaluación se repetirá sobre múltiples instantes y condiciones de
circulación. Por tanto, además de los valores agregados de MAE y RMSE, deberá
examinarse la dispersión de los errores para identificar si el desempeño se
mantiene estable o se deteriora en determinados estados de la red. En esta
etapa el objetivo es caracterizar el desempeño del componente de estimación,
no determinar todavía el efecto final sobre las rutas generadas por el
DVRPTW.

La evaluación operacional completa corresponderá a la integración posterior
con el modelo de optimización. En esa etapa podrán analizarse, entre otros
aspectos, la factibilidad de las rutas, la calidad de la solución, el costo
resultante y el tiempo total de respuesta. Mantener separados ambos niveles de
evaluación permitirá determinar primero si el mecanismo de consultas guía
produce una representación suficientemente precisa y rápida de los tiempos de
viaje y, posteriormente, si dicha representación se traduce en mejores
decisiones de planeación o replaneación.

La mejora predictiva no se interpretará por sí sola como una mejora integral
del sistema. Una vez incorporadas las estimaciones al modelo de ruteo, la
evaluación deberá considerar también efectos sobre los actores involucrados y
sobre criterios transdisciplinarios relevantes, tales como cumplimiento de
ventanas de servicio, tiempo de operación de los conductores, utilización de
recursos, costos y estabilidad de la operación. De esta forma, las métricas de
error de la capa de estimación y las métricas operacionales del modelo de ruteo
se mantendrán diferenciadas, pero vinculadas dentro de la evaluación final.

\section{Alcance del avance y etapas posteriores}

El presente avance deja establecida la formulación necesaria para estudiar la
actualización de los tiempos de viaje utilizados por el modelo de ruteo. Se
definieron el punto de integración con el VRPTW/DVRPTW de referencia, el papel
de PEMS-BAY como entorno experimental, la construcción de tiempos de
referencia, el esquema de consultas guía y los criterios con los que podrá
evaluarse posteriormente la precisión y rapidez del componente de estimación.

El alcance actual no comprende todavía la selección definitiva ni la
implementación comparativa de un modelo de inteligencia artificial, ni la
evaluación de su efecto sobre las rutas generadas por el optimizador. La
revisión metodológica permite identificar enfoques compatibles con la
formulación propuesta, pero su implementación y comparación experimental
corresponderán a etapas posteriores del proyecto.

En los siguientes avances se implementará o reproducirá primero el método de
referencia para la optimización y recuperación de rutas. Posteriormente podrán
implementarse y compararse alternativas para la estimación de tiempos de viaje
utilizando los escenarios y métricas definidos en este documento. Una vez
seleccionada una alternativa suficientemente precisa y rápida, sus estimaciones
\(\widehat{\tau}_{ij}(t)\) podrán incorporarse a
\(\widehat{T}^{\,\mathrm{rec}}\) para analizar su efecto sobre la factibilidad,
calidad y tiempo de respuesta del proceso de replaneación.

La predicción de condiciones futuras para la planeación previa permanece como
uno de los dos componentes centrales de la arquitectura propuesta, aunque su
validación experimental detallada se desarrollará en una etapa posterior a la
primera prueba de concepto enfocada en actualización rápida. Asimismo, la
incorporación de eventos contextuales mediante representaciones semánticas, la
validación con información operativa de una empresa, la utilización de
servicios externos de tráfico en tiempo real y el desarrollo de estrategias
más avanzadas para seleccionar las consultas guía se consideran extensiones
posteriores de la implementación experimental.

\subsection{Conclusión de preparación de datos en CRISP-ML}

La fase de preparación de datos queda delimitada por la conversión de las
mediciones de movilidad en tiempos de referencia, la identificación de
entradas disponibles durante la operación y los criterios de construcción,
escalamiento y selección o extracción de características. La separación
entre entradas y objetivos, junto con el ajuste del procesamiento únicamente
sobre entrenamiento, busca reducir el riesgo de fuga de información y de
una evaluación excesivamente favorable.

Para concluir esta fase será necesario ejecutar el procesamiento y documentar
la calidad de los datos, las distribuciones antes y después de las
transformaciones, las características conservadas y la reducción obtenida.
El diseño aquí presentado no permite afirmar todavía mejoras de rendimiento
ni de convergencia. En CRISP-ML, el paso al modelado deberá apoyarse en un
conjunto preparado de forma reproducible y en evidencia de que las decisiones
adoptadas son compatibles con la información disponible durante la operación.

% Bibliografía integrada: mismos trabajos citados y datos bibliográficos.
\begin{thebibliography}{99}
\bibitem[Berth \& Trigos(2025)]{berth_trigos_2025}
Berth, M. O., \& Trigos, F. (2025).
A Transdisciplinary Framework for Handling Disruptions in Vehicle Routing with Customer Time Windows. En F. Trigos et al. (Eds.),
\textit{Transdisciplinarity for a Better World} (pp. 201--210).
\url{https://doi.org/10.3233/ATDE251092}

\bibitem[Berth(2025)]{berth_thesis_2025}
Berth, M. O. (2025).
\textit{Transdisciplinary Approach to Disruptions in Vehicle Routing Problems} [Tesis de maestría]. Reutlingen University.
Written at Instituto Tecnológico y de Estudios Superiores de Monterrey.

\bibitem[California Department of Transportation(s.f.)]{caltrans_pems}
California Department of Transportation (s.f.).
\textit{Performance Measurement System ({PeMS}) Data Source}. Consultado el 9 de octubre de 2026.
\url{https://dot.ca.gov/programs/traffic-operations/mpr/pems-source}

\bibitem[Cebecauer et al.(2018)]{cebecauer2018integrated}
Cebecauer, M., Jenelius, E., \& Burghout, W. (2018).
Integrated Framework for Real-Time Urban Network Travel Time Prediction on Sparse Probe Data.
\textit{IET Intelligent Transport Systems, 12}(1), 66--74.
\url{https://doi.org/10.1049/iet-its.2017.0113}

\bibitem[Du et al.(2025)]{du2025survey}
Du, K., Guo, X., Li, L., Song, J., Shi, Q., Hu, M., \& Fang, J. (2025).
Traffic Prediction in Time Series, Spatialtemporal, and {OD} Data: A Systematic Survey.
\textit{Journal of Traffic and Transportation Engineering (English Edition), 12}(3), 666--700.
\url{https://doi.org/10.1016/j.jtte.2025.03.001}

\bibitem[Gentile et al.(2020)]{gentile2020fast}
Gentile, G., Meschini, L., Tiddi, D., Fiorani, A., \& Attanasi, A. (2020).
Fast Estimation of Point-to-Point Travel Times for Real-Time Vehicle Routing.
\textit{Transportation Research Procedia, 47}, 401--408.
\url{https://doi.org/10.1016/j.trpro.2020.03.115}

\bibitem[Hofleitner et al.(2012)]{hofleitner2012learning}
Hofleitner, A., Herring, R., Abbeel, P., \& Bayen, A. M. (2012).
Learning the Dynamics of Arterial Traffic From Probe Data Using a Dynamic Bayesian Network.
\textit{IEEE Transactions on Intelligent Transportation Systems, 13}(4), 1679--1693.
\url{https://doi.org/10.1109/TITS.2012.2200474}

\bibitem[Jenelius \& Koutsopoulos(2018)]{jenelius2018urban}
Jenelius, E., \& Koutsopoulos, H. N. (2018).
Urban Network Travel Time Prediction Based on a Probabilistic Principal Component Analysis Model of Probe Data.
\textit{IEEE Transactions on Intelligent Transportation Systems, 19}(2), 436--445.
\url{https://doi.org/10.1109/TITS.2017.2703652}

\bibitem[Lei et al.(2024)]{lei2024conditional}
Lei, D., Xu, M., \& Wang, S. (2024).
A Conditional Diffusion Model for Probabilistic Estimation of Traffic States at Sensor-Free Locations.
\textit{Transportation Research Part C: Emerging Technologies, 166}, 104798.
\url{https://doi.org/10.1016/j.trc.2024.104798}

\bibitem[Li et al.(2018)]{li2018dcrnn}
Li, Y., Yu, R., Shahabi, C., \& Liu, Y. (2018).
\textit{Diffusion Convolutional Recurrent Neural Network: Data-Driven Traffic Forecasting}. International Conference on Learning Representations.

\bibitem[Shah \& Piantanakulchai(2026)]{shah2026active}
Shah, S. M. T., \& Piantanakulchai, M. (2026).
Ensemble Machine Learning for Travel Time Prediction: An Active Learning-Assisted Approach.
\textit{Results in Engineering, 29}, 108712.
\url{https://doi.org/10.1016/j.rineng.2025.108712}

\bibitem[Wang et al.(2025)]{wang2025od}
Wang, R., Guo, Q., Dai, S., Deng, L., Xiao, Y., \& Jia, C. (2025).
An {OD} Time Prediction Model Based on Adaptive Graph Embedding.
\textit{Physica A: Statistical Mechanics and its Applications, 657}, 130217.
\url{https://doi.org/10.1016/j.physa.2024.130217}

\bibitem[Zhan et al.(2016)]{zhan2016bayesian}
Zhan, X., Ukkusuri, S. V., \& Yang, C. (2016).
A Bayesian Mixture Model for Short-Term Average Link Travel Time Estimation Using Large-Scale Limited Information Trip-Based Data.
\textit{Automation in Construction, 72}, 237--246.
\url{https://doi.org/10.1016/j.autcon.2015.12.007}

\bibitem[Zhou et al.(2025)]{zhou2025mogernn}
Zhou, Q., Zhang, Y., Makridis, M. A., Kouvelas, A., Wang, Y., \& Hu, S. (2025).
{MoGERNN}: An Inductive Traffic Predictor for Unobserved Locations.
\textit{Transportation Research Part C: Emerging Technologies, 174}, 105080.
\url{https://doi.org/10.1016/j.trc.2025.105080}
\end{thebibliography}
\end{document}